%%%%%%%%%%%%%%%%%%%%%%%%%%%%%%%%%%
%%%%% ----- New colors ----- %%%%%
%%%%%%%%%%%%%%%%%%%%%%%%%%%%%%%%%%

\usepackage{tcolorbox}                                  % ----- Boxes environment
    \tcbuselibrary{skins,listings,listingsutf8,breakable}             % ----- All libraries
\usepackage[bold=0.7]{xfakebold}


%%%%%%%%%%%%%%%%%%%%%%%%%%%%%%%%%%
%%%%% ----- New colors ----- %%%%%
%%%%%%%%%%%%%%%%%%%%%%%%%%%%%%%%%%

\definecolor{persianblue}{HTML}{221177} % ----- New blues
\definecolor{lightblue}{HTML}{EEFFFF}
\definecolor{softblue}{HTML}{6655BB}

\definecolor{maroon}{HTML}{992222}      % ----- New reds
\definecolor{lightred}{HTML}{FAEEEE}

\definecolor{calgreen}{HTML}{005500}    % ----- New greens
\definecolor{lightgreen}{HTML}{EEFFEE}
\definecolor{softgreen}{HTML}{449944}

\definecolor{brickorange}{HTML}{CC7033} % ----- New oranges
\definecolor{lightorange}{HTML}{FFF2E6}

\definecolor{newred}{HTML}{CC0000}     % ----- Colors for geometry
\definecolor{newblue}{HTML}{0000CC}
\definecolor{newgreen}{HTML}{00AA00}
\definecolor{neworange}{HTML}{FFAA00}





%%%%%%%%%%%%%%%%%%%%%%%%%%%%%%%%%
%%%%% ----- \theoreme ----- %%%%%
%%%%%%%%%%%%%%%%%%%%%%%%%%%%%%%%%

\newtcolorbox{theorembox}[2][]{% ----- Theorem box
enhanced,
arc=5pt,
frame style = {left color=persianblue , right color=softblue},
colback=lightblue,
fonttitle=\bfseries,
lifted shadow={1mm}{-1.5mm}{2mm}{0.3mm},
hbox,
title={#2},
attach boxed title to top left={xshift=5mm,yshift=-2mm},
boxed title style = {colback = persianblue , colframe = persianblue},
#1}


\newcommand\theoreme[2]{% ----- Theorem command
    \begin{theorembox}[]{\faBook~~#1~~\faBook}
        \begin{varwidth}{\linewidth-15pt }
        \restoreskip \vspace{0pt}
        #2
        \end{varwidth}
    \end{theorembox}
}



%%%%%%%%%%%%%%%%%%%%%%%%%%%%%%%%%%%
%%%%% ----- \definition ----- %%%%%
%%%%%%%%%%%%%%%%%%%%%%%%%%%%%%%%%%%

\newtcolorbox{definitionbox}[2][]{% ----- Definition box
enhanced,
arc=5pt,
frame style = {left color=calgreen , right color=softgreen},
colback=lightgreen,
fonttitle=\bfseries,
lifted shadow={1mm}{-1.5mm}{2mm}{0.3mm},
hbox,
title={#2},
attach boxed title to top left={xshift=5mm,yshift=-2mm},
boxed title style = {colback = calgreen , colframe =calgreen},
#1}


\newcommand\definition[2]{% ----- Definition command
    \begin{definitionbox}[]{\faUniversity~~#1~~\faUniversity}
        \begin{varwidth}{\linewidth-15pt}
        \restoreskip \vspace{0pt}
        #2
        \end{varwidth}
    \end{definitionbox}
}



%%%%%%%%%%%%%%%%%%%%%%%%%%%%%%%
%%%%% ----- \astuce ----- %%%%%
%%%%%%%%%%%%%%%%%%%%%%%%%%%%%%%

\newtcolorbox{astucebox}[1][]{% ----- astuce box
enhanced,
arc=10pt,
frame hidden,
borderline={1mm}{0mm}{maroon,dotted},
colback=lightred,
fonttitle=\bfseries,
hbox,
title={\faLightbulb~~Astuce~~\faLightbulb},
attach boxed title to top left={xshift=5mm,yshift=-2mm},
boxed title style = {colback = maroon , colframe =maroon},
#1}


\newcommand\astuce[1]{% ----- astuce command
    \begin{astucebox}[]
        \begin{varwidth}{\linewidth-15pt }
        \restoreskip \vspace{0pt}
        #1
        \end{varwidth}
    \end{astucebox}
}


%%%%%%%%%%%%%%%%%%%%%%%%%%%%%%%
%%%%% ----- \remarque ----- %%%%%
%%%%%%%%%%%%%%%%%%%%%%%%%%%%%%%

\newtcolorbox{remarquebox}[1][]{% ----- remarque box
enhanced,
arc=10pt,
frame hidden,
borderline={1mm}{0mm}{maroon,dotted},
colback=lightred,
fonttitle=\bfseries,
hbox,
title={\faExclamation~~Remarque~~\faExclamation},
attach boxed title to top left={xshift=5mm,yshift=-2mm},
boxed title style = {colback = maroon , colframe =maroon},
#1}


\newcommand\remarque[1]{% ----- remarque command
    \begin{remarquebox}[]
        \begin{varwidth}{\linewidth-15pt }
        \restoreskip \vspace{0pt}
        #1
        \end{varwidth}
    \end{remarquebox}
}



%%%%%%%%%%%%%%%%%%%%%%%%%%%%%%%%
%%%%% ----- \methode ----- %%%%%
%%%%%%%%%%%%%%%%%%%%%%%%%%%%%%%%

\newtcolorbox{methodebox}[2][]{% ----- methode box
enhanced,
arc=10pt,
frame hidden,
borderline={1mm}{0mm}{maroon,dotted},
colback=lightred,
fonttitle=\bfseries,
hbox,
title={#2},
attach boxed title to top left={xshift=5mm,yshift=-2mm},
boxed title style = {colback = maroon , colframe =maroon},
#1}


\newcommand\methode[2]{% ----- methode command
    \begin{methodebox}[]{\faCogs~~#1~~\faCogs}
        \begin{varwidth}{\linewidth-20pt }
        \restoreskip \vspace{0pt}
        #2
        \end{varwidth}
    \end{methodebox}
}



%%%%%%%%%%%%%%%%%%%%%%%%%%%%%%%%
%%%%% ----- \exemple ----- %%%%%
%%%%%%%%%%%%%%%%%%%%%%%%%%%%%%%%

\newtcolorbox{exemplebox}[1][]{% ----- exemple box
enhanced, breakable,
arc = 0pt ,
frame hidden,
borderline west={2pt}{0pt}{brickorange},
interior style = {top color = lightorange , bottom color = lightorange!15!white},
fonttitle=\bfseries,
minipage,
#1}


\newcommand\exemple[2]{% ----- exemple command
    \begin{exemplebox}[]
        \textbf{{\color{brickorange}Exemple :} #1} \\[5pt]
        \restoreskip \vspace{0pt}
        #2
    \end{exemplebox}
    {\color{brickorange}
    \vspace{-2pt}
    \hrule height 2pt depth 0pt width 3cm
    \vspace{-11pt} \hspace{2.85cm} \setBold{$\square$}}
}




%%%%%%%%%%%%%%%%%%%%%%%%%%%%%%%%%%%%%%
%%%%% ----- \demonstration ----- %%%%%
%%%%%%%%%%%%%%%%%%%%%%%%%%%%%%%%%%%%%%

\newtcolorbox{demonstrationbox}[1][]{% ----- demonstration box
enhanced, breakable,
arc = 0pt
frame hidden,
borderline west={2pt}{0pt}{VioletRed3},
interior style = {top color = Thistle1 , bottom color = Thistle1!15!white},
fonttitle=\bfseries,
hbox,
#1}


\newcommand\demonstration[1]{% ----- demonstration command
    \begin{demonstrationbox}[]
        \begin{varwidth}{\linewidth-15pt }
        {\vspace{-5pt}\textbf{\color{VioletRed3}Démonstration :}}\\[5pt]
        \restoreskip \vspace{0pt}
        #1
        \end{varwidth}
    \end{demonstrationbox}
    \vspace{-8 pt}
    {\color{VioletRed3}
    \vspace{-2pt}
    \hrule height 2pt depth 0pt width 3cm
    \vspace{-11pt} \hspace{2.85cm} \setBold{$\square$}}
}



%%%%%%%%%%%%%%%%%%%%%%%%%%%%%%%%%
%%%%% ----- \exercice ----- %%%%%
%%%%%%%%%%%%%%%%%%%%%%%%%%%%%%%%%

\newtcolorbox[auto counter , number within=section , number freestyle={\noexpand\arabic{\tcbcounter}} ]{exercicebox}[2][]{% ----- Exercice box
enhanced,
arc=1pt,
colframe=black!80,
colback=black!10,
fonttitle=\bfseries,
hbox,
title={Exercice n° \thetcbcounter \, : #2 },
attach boxed title to top*,
boxed title style = {colback = black!80},
#1}


\newcommand\exercice[2]{% ----- Exercice command
    \begin{exercicebox}[]{#1}
        \begin{varwidth}{\linewidth-30pt} 
         #2
        \end{varwidth}
    \end{exercicebox}
}



%%%%%%%%%%%%%%%%%%%%%%%%%%%%%%%%
%%%%% ----- \calculs ----- %%%%%
%%%%%%%%%%%%%%%%%%%%%%%%%%%%%%%%

\newtcolorbox{calculbox}[1][]{% ----- Calculs box
enhanced,
arc=1pt,
colbacktitle=black!25!white,
coltitle=black,
colback=white,
colframe=black!25!white,
hbox,
title={Aides de calcul},
attach boxed title to top*,
#1}


\newcommand\calculs[1]{% ----- Calculs command
    \begin{calculbox}
        \restoreskip
        #1
    \end{calculbox}
}


